\section{Related work}\label{sec:related}

\subsection{Register projects define opportunities through specific releases}

The Dutch 2018 file formalised population-wide linkage across administrative sources. Its data brief defines family, household, work, school and neighbour ties and sets out the opportunities and coverage limitations of a graph constructed this way \citep{vanderlaan2022,vanderlaan2023}. The Dutch longitudinal files are a distinct data product. Their methodology changed the neighbour, work, household, education and family layers together \citep{cbs2023,das2023}. For comparisons, the release name and date determine which candidate relationships and nomination constraints are in force. Treating the annual files as repeated draws from the original specification would attribute design changes to social change.

The Danish project builds annually resolved networks and makes a person-to-group representation available under controlled access. The published one-mode network and the unprojected memberships support different questions. A researcher with the group data can ask how an alternative projection changes a statistic. A researcher with the projected graph alone generally cannot recover all unsampled co-members \citep{cremers2025}. The Swedish project uses different administrative units for property, education and work, and combines information across years for parts of the family construction \citep{panayiotou2025,wohlert2024}. The release-specific scripts provide a valuable independent source about how the published description was executed, including the link-function constants and cohort handling \citep{secode2026}. Those scripts and a card can therefore agree on a particular operator while having different input contracts or eligibility coverage.

\subsection{Applications make release semantics consequential}

The Dutch anatomy study characterises degree, multiplex ties and closure across domains \citep{bokanyi2023}. The income segregation study measures exposure in a network extending beyond residential boundaries, and Statistics Netherlands uses the yearly files for segregation by origin \citep{kazmina2024,vanklingeren2024}. A separate longitudinal preprint examines changes in cohesion across twelve annual snapshots \citep{bokanyi2026}. These applications make a coding difference concrete. A neighbour definition affects the set of candidate alters for exposure. The treatment of a household-and-family pair affects overlap, and the ratio of closed to eligible alter pairs depends on how repeated ties are counted.

Other uses link constructed network summaries to measures collected independently. \citet{pitoski2025} combine Dutch register and survey information to study social cohesion. \citet{menyhert2025} compare a register network with online social ties at municipality level and report similarities in aggregate connectivity. Municipal agreement and validation of individual dyads are different targets. The latter would require identified edges from a contact or relationship measure with a matching unit and reference period. In kinship research, \citet{debel2025} replicate a Swedish demographic construction on Dutch registers and compare explicitly defined kin categories. Their design illustrates the value of a common relationship definition when register coverage differs.

The work layer has an analogous history. \citet{lengyel2017} varied tie-selection rules and caps in Swedish matched employer--employee records. Their aim was to model co-worker networks for regional labour mobility and productivity. Studies of worker sorting, low-wage transitions and spatial colleague connectivity likewise use linked employment records to construct different measures of social opportunity \citep{eliason2023,baranowska2023,song2026}. A work tie may represent previous co-employment, an opportunity within a current site, or proximity among colleagues' homes. Code that exposes the workplace unit and nomination rule makes these design choices visible before such measures are compared.

\subsection{Projection introduces structure before network analysis begins}

A person-to-group table is naturally a bipartite graph. A full projection places a clique on every group. A large group contributes quadratically many unordered candidate pairs, so group sizes can dominate a degree or triangle calculation \citep{latapy2008}. Random bipartite and projected graph models formalise how degree and clustering can arise from affiliations \citep{newman2001}. \citet{brunson2015} defines closure on affiliations to address the influence of repeat attendance. These results motivate viewing group definitions, overlap and projection as separate inputs. A school year, a classroom and a graduation cohort can share a data source while producing distinct co-member sets.

Limited nomination has close precedents in fixed-choice data. \citet{kossinets2006} studies missing actors, missing affiliations and fixed-choice censoring. \citet{harling2018} examine degree truncation in network-based transmission models. \citet{griffith2022} studies censoring in peer effect estimation, and \citet{neal2024} examines whether the order of a roster affects who is nominated. Register projections have a different data-generating sequence. The cap is a stated transformation of an administrative affiliation, while a fixed-choice survey cap limits the report of an underlying relationship. Both settings make it necessary to distinguish the degree a person can nominate from the degree an analyst obtains after nominations are united.

\subsection{Multiplex statistics retain the semantics of repeated relationships}

A binary union gives each pair equal weight. A multiplex representation retains the counts and identities of the contributing layers. \citet{battiston2014} use layer degrees to define participation and other structural measures. \citet{cozzo2015} give several ways to generalise triangle closure to multiplex networks. Their treatment makes clear that counting which layer closes a triplet and choosing the denominator are separate steps. \citet{barrat2004} likewise distinguish weighted from unweighted clustering by keeping the allocation of weight around an ego in the definition. In a register graph, household and kinship can describe the same people. Making the closure convention an explicit argument therefore makes repeated analyses comparable at the level of their estimands.

Construction comparisons also differ from inferential tests on a fixed network. A changed network can change both a descriptor and its dependence structure. In animal network analysis, \citet{farine2022} and \citet{hart2022} study cases where permutation tests fail to represent the relevant null or control dependence. That work motivates treating the software's label shuffle as a particular diagnostic. It gives no automatic licence to interpret a nominal node-level $p$-value when the data are a projection of shared institutions.

\subsection{Synthetic populations and graph systems solve adjacent problems}

Synthetic population software builds plausible agents, groups and contact networks for particular analysis goals. \citet{jiang2024} assemble geographically explicit US synthetic populations and stylised home, school and work networks. \citet{tulchinsky2026} model geographically detailed contacts from public inputs. \citet{chen2025} describe a scalable epidemic platform with dynamic network and intervention machinery. These workflows connect a population representation to a contact or transmission model. The question here concerns what happens between membership data and an administrative graph summary. It can be investigated on the same inputs even when no transmission process is defined.

MARS uses spatial affiliations to generate register-like multiplex networks \citep{hamilton2026}. Its reported clique model provides a baseline for studying how affiliation geography shapes closure. The comparison supported by \builtnet{} also allows the same memberships to be projected through capped, spatial and random alternatives. Threadle is closer in setting. It provides storage, query and processing primitives for large register-derived multimode networks, together with guidance for Swedish register workflows \citep{nordlundthreadle2026,nordlund2026}. General software for multilayer and higher-order networks provides representations and analyses after the network definition has been chosen. \texttt{pymnet} supports multilayer graphs, and XGI and \texttt{hypergraphx} support explicit group-interaction representations \citep{nurmi2024,landry2023,lotito2023}. Such representations can preserve group information discarded by a simple graph. The construction-card interface addresses the choice of which graph to produce from that information.

\subsection{Executable provenance makes the comparison repeatable}

FAIR principles for research software emphasise identifiable versions, provenance and reuse \citep{barker2022}. RO-Crate allows research objects and their metadata to be packaged together \citep{soiland2022}. A construction card records a narrower unit of provenance. It identifies the source of the fields that determine a projection and connects those fields to an executable rule set. This permits a direct test that an analysis changes a cap, a radius or an eligibility rule as intended. In an adjacent network measurement problem, \citet{luppi2024} compared functional-connectivity pipelines across reconstruction choices. The analogous design here is to cross the construction cards with a fixed administrative input while treating each shared output definition as part of the comparison.

Table~\ref{tab:prior} states the contribution against these precedents. Among the compared artefacts, \builtnet{} combines a single membership interface with a catalogue of source-attributed rules spanning Dutch, Danish and Swedish releases, including executable alternatives for ambiguous fields. This claim is testable by inspecting the card inheritance, changing one card field, and building both alternatives on the same membership rows. It does not depend on priority for varying a cap, writing national construction code or storing a large network.
% Generated by src/make_tables.py.
\begin{table}[!htbp]
\centering\small
\caption{Comparison of prior work with the software contribution. Each row describes the cited artefact.}
\label{tab:prior}
\begin{tabular}{@{} >{\raggedright\arraybackslash}p{0.27\linewidth} >{\raggedright\arraybackslash}p{0.29\linewidth} >{\raggedright\arraybackslash}p{0.34\linewidth} @{}}
\toprule
\rowcolor{tableblue} \textbf{Work} & \textbf{Main reusable object or comparison} & \textbf{Construction choices exposed} \\
\midrule
Dutch and Danish releases \citep{vanderlaan2023,cbs2023,cremers2025} & Specific network definitions; Denmark also provides two-mode data & Release documentation describes layer rules and eligibility \\
Swedish release \citep{panayiotou2025,secode2026} & Source-specific construction code and network description & Executable steps for the Swedish construction \\
Co-worker sensitivity \citep{lengyel2017} & Alternative tie draws on Swedish employment records & Varies the cap and choice of colleague within a work application \\
MARS \citep{hamilton2026} & Spatial generator for register-like multilayer graphs & Uses group assignment and clique projection in the reported model \\
Threadle \citep{nordlundthreadle2026,nordlund2026} & Large register-network storage, processing and queries & Supports register workflows with multimode or multilayer relations \\
General graph software \citep{nurmi2024,landry2023,lotito2023} & Analytic representations for multilayer and group networks & Data structures and algorithms for their respective network classes \\
\textit{builtnet} \citep{builtnet2026} & Executable cards and common membership input & Source-tagged options across documented Dutch, Danish and Swedish constructions \\
\bottomrule
\end{tabular}
\end{table}

\FloatBarrier
